%%%%%%%%%%%%%%%%%%%%%%%%%%%%%%%%%%%%%%%%%%%%%%%%%%%%%%%%%%%%%
% ==================== 5.4：问题四的模型建立与求解 ====================
% ✓ 自查：综合使用前三问成果；最终结论逐项回应题目；检验服务于结论而非形式完整。

\subsection{问题四的模型建立与求解}

\subsubsection{建模思路与模型选择}

问题四要求综合前三问结果完成【最终预测、决策、优化或机理解释任务】。
本问以【问题一结果】、【问题二结果】和【问题三结果】作为【输入、参数或备选方案】，
同时考虑【最终新增条件或现实限制】，需要回答【核心问题】。
据此选择【主模型名称或综合分析方法】完成【最终任务】，并以【评价标准或基准方案】衡量结果质量。
如有必要的数据整理，只说明与最终模型直接相关的【处理内容与理由】。

\subsubsection{模型的建立}

定义【最终决策变量、输出变量或状态变量】为【符号及含义】，定义【新增参数】为【符号及含义】。
结合前三问结论和【现实约束】，建立【核心方程、目标函数、决策规则或机理关系】：

\begin{equation}
  \text{【问题四核心方程、目标函数、决策规则或机理关系】},
  \label{eq:q4_core}
\end{equation}

\begin{equation}
  \text{【问题四约束、边界条件或适用条件】}.
  \label{eq:q4_constraint}
\end{equation}

\noindent\textbf{模型汇总：}
综合全部输入、目标、约束和变量范围，得到问题四的完整模型

\begin{equation}
  \mathcal{M}_4:\quad
  \left\{
  \begin{aligned}
    &\text{【最终模型表达式】},\\
    &\text{【全部主要约束或决策规则】},\\
    &\text{【变量范围、初始条件或边界条件】}.
  \end{aligned}
  \right.
  \label{eq:q4_summary}
\end{equation}

\subsubsection{模型的求解}

\noindent\textbf{Step 1：}汇总前三问的【真实输出】和问题四的【新增输入】，统一符号、单位与评价口径。

\noindent\textbf{Step 2：}采用【算法、求解器或综合分析流程】计算【最终预测、决策变量或候选方案】。

\noindent\textbf{Step 3：}根据【评价标准、可行性条件或决策准则】筛选最终结果，并形成【题目要求的输出形式】。

求解使用【软件、求解器或程序语言】完成，关键参数与终止条件为【真实设置及依据】，完整程序见附录【编号】。

\subsubsection{求解结果与分析}

最终得到【真实的预测值、最优方案、决策建议或机理结论】。
由【图或表编号】可知，【主要规律或方案特征】，其产生原因是【结合模型与现实条件解释原因】。
上述结果分别回答了【题目要求一】与【题目要求二】，最终结论为【完整回答问题四】。

\subsubsection{模型检验与灵敏度分析}

围绕最终结论，选择【与主模型匹配的检验方法】验证【结果精度、约束可行性、稳定性或机理合理性】。
检验得到【真实指标、对照结果或边界表现】，依据【判断标准】可得【检验结论】。
若存在会改变最终决策的【现实不确定参数】，则在【有依据的范围】内分析其影响，
说明【主要趋势、稳定区间或策略切换点】；不得用程序运行正常代替数学意义上的模型检验。
